\chapter{Réalisation}
\label{chap-realisation}

\section*{Introduction}
\addcontentsline{toc}{section}{Introduction}
Ce chapitre décrit la mise en œuvre de la plateforme de gestion des projets de stage, à partir des besoins et des choix de conception définis précédemment. Il présente les technologies utilisées, l'architecture applicative, les fonctionnalités disponibles pour chaque rôle, la chaîne DevSecOps et les interfaces réalisées.

\section{Choix technologiques}
\begin{table}[h]
\centering
\caption{Technologies principales}
\label{tab-technologies}
\begin{tabularx}{\textwidth}{p{3.2cm}Y}
\toprule
\textbf{Technologie} & \textbf{Utilisation dans le projet} \\
\midrule
Spring Boot & Exposition de l'API REST, logique métier, sécurité et accès aux données. \\
React + TypeScript & Interface web par rôle, tableaux de bord, gestion des projets et des tâches. \\
PostgreSQL & Persistance des utilisateurs, projets, tâches, soumissions, pièces jointes et commentaires. \\
Jenkins & Automatisation de la construction, des tests, des analyses et du déploiement. \\
Docker & Construction et exécution isolée des services frontend et backend. \\
SonarQube & Analyse de la qualité et de la maintenabilité du code, avec Quality Gate. \\
GitLeaks / OWASP / Trivy & Détection de secrets, analyse des dépendances et scan des images de conteneurs. \\
\bottomrule
\end{tabularx}
\end{table}

Le Tableau~\ref{tab-technologies}, intitulé \textit{Technologies principales}, résume les outils retenus pour la réalisation. Il distingue les composants de l'application, les outils d'automatisation et les contrôles de qualité et de sécurité intégrés à la chaîne DevSecOps.

\section{Architecture applicative}
L'application est structurée en deux parties. Le frontend React fournit des espaces adaptés au SuperAdmin, au Supervisor et à l'Intern. Le backend Spring Boot applique les règles de gestion, assure l'authentification et expose les ressources via une API REST. PostgreSQL garantit la persistance des données métier.

Les échanges sont sécurisés grâce à une authentification par jeton JWT et à une autorisation fondée sur les rôles. La séparation des responsabilités facilite l'évolution indépendante des interfaces et des services métier.

\begin{figure}[h]
\centering
\begin{tikzpicture}[font=\small, box/.style={draw,rounded corners,fill=lightblue,minimum width=2.4cm,minimum height=.8cm,align=center}, >=Latex]
\node[box] (client) {Navigateur\\React + TypeScript};
\node[box,right=1.2cm of client] (api) {API REST\\Spring Boot};
\node[box,right=1.2cm of api] (db) {PostgreSQL};
\node[box,below=1.3cm of api] (ai) {Service IA\\Génération de brouillons};
\node[box,above=1.3cm of api] (jenkins) {Jenkins\\CI/CD et DevSecOps};
\draw[-Latex,thick] (client)--node[above]{HTTPS / API}(api);
\draw[-Latex,thick] (api)--(db);
\draw[-Latex,thick] (api)--(ai);
\draw[-Latex,thick] (jenkins)--(api);
\end{tikzpicture}
\caption{Architecture du projet de la plateforme}
\label{fig-architecture-plateforme}
\end{figure}

La Figure~\ref{fig-architecture-plateforme}, intitulée \textit{Architecture du projet de la plateforme}, illustre la communication entre le navigateur React, l'API Spring Boot et PostgreSQL. Elle met également en valeur le service IA, utilisé pour générer des brouillons, ainsi que Jenkins, qui automatise la chaîne CI/CD et les contrôles DevSecOps.

\section{Fonctionnalités réalisées}
\subsection{Espace SuperAdmin}
Le SuperAdmin gère les comptes des encadrants et des stagiaires, crée et affecte les projets, et consulte une vue globale des statuts. Il peut également soumettre une instruction au service IA afin d'obtenir un brouillon de projet et de tâches. Ce contenu n'est enregistré qu'après relecture, modification éventuelle et confirmation explicite.

\subsection{Espace Supervisor}
Le Supervisor accède aux projets et stagiaires qui lui sont affectés. Il crée les tâches, précise leur priorité et leur échéance, suit leur état, consulte les soumissions et formule des commentaires.

\subsection{Espace Intern}
L'Intern consulte son projet et ses tâches, soumet ses réponses et associe des fichiers ou des liens externes. Il peut ensuite consulter les commentaires reçus et suivre l'avancement de son travail.

\section{Chaîne DevSecOps}
La chaîne Jenkins automatise les étapes suivantes :
\begin{enumerate}
  \item récupération du code source ;
  \item recherche de secrets avec GitLeaks ;
  \item construction et tests du backend Maven et du frontend React ;
  \item analyse des dépendances Java avec OWASP Dependency-Check et des dépendances frontend avec \texttt{pnpm audit} ;
  \item analyse SonarQube suivie du contrôle du Quality Gate ;
  \item construction des images Docker ;
  \item analyse de sécurité des images avec Trivy ;
  \item déploiement et vérification de l'état des conteneurs.
\end{enumerate}
Cette approche déplace les vérifications de qualité et de sécurité au plus près du cycle de développement, afin de détecter les problèmes avant la mise à disposition de l'application.

\section{Interfaces réalisées}
Les interfaces suivantes illustrent les principales fonctionnalités livrées. Elles sont organisées par parcours utilisateur afin de distinguer la page d'authentification, l'espace SuperAdmin, l'espace Supervisor et l'espace Intern.

\newcommand{\screenpage}[5]{%
  \begin{center}
    \includegraphics[width=#2]{#1}
    \captionof{figure}{#3}
    \label{#4}
  \end{center}
  \noindent\textbf{Explication de la Figure~\ref{#4} -- \textit{#3}.} \newline
  #5
  \medskip
}

\newcommand{\screenpair}[6]{%
  \begin{center}
    \begin{minipage}[t]{0.48\textwidth}
      \centering
      \includegraphics[width=\textwidth]{#1}
      \captionof{figure}{#2}
      \label{#3}
    \end{minipage}\hfill
    \begin{minipage}[t]{0.48\textwidth}
      \centering
      \includegraphics[width=\textwidth]{#4}
      \captionof{figure}{#5}
      \label{#6}
    \end{minipage}
  \end{center}
  \medskip
}

\clearpage
\subsection{Authentification}
\screenpage{login-page.png}{\textwidth}{Page de connexion de la plateforme}{fig-login-page}{%
La page de connexion constitue le point d'entrée de la plateforme. Elle permet à l'utilisateur de s'authentifier avec son adresse e-mail et son mot de passe. L'interface met aussi en évidence des fonctionnalités prévues en évolution, notamment la mémorisation de session, l'authentification SSO et l'identification par badge.
}

\subsection{Interfaces de l'espace SuperAdmin}
\screenpage{superadmin-dashboard.png}{\textwidth}{Tableau de bord de l'espace SuperAdmin}{fig-superadmin-dashboard}{%
Après authentification, le SuperAdmin accède à un tableau de bord de synthèse. Cette vue présente les indicateurs principaux de la plateforme, tels que le nombre de projets actifs, le nombre de membres et le volume de tâches terminées. Une zone d'activité récente facilite le suivi des dernières actions réalisées par les utilisateurs.
}

\screenpage{superadmin-projects.png}{\textwidth}{Liste des projets dans l'espace SuperAdmin}{fig-superadmin-projects}{%
La page de gestion des projets permet au SuperAdmin de consulter l'ensemble des projets, leur statut et l'encadrant affecté. Depuis cette interface, il peut créer un nouveau projet, modifier les informations d'un projet existant ou supprimer une entrée si nécessaire.
}

\screenpage{superadmin-team.png}{\textwidth}{Gestion des comptes utilisateurs}{fig-superadmin-team}{%
L'interface de gestion d'équipe centralise les comptes de la plateforme. Elle affiche le nom, l'adresse e-mail, le rôle et le statut de chaque utilisateur. Le SuperAdmin peut ainsi administrer les comptes SuperAdmin, Supervisor et Intern depuis une vue unique.
}

\screenpage{superadmin-register-user.png}{0.82\textwidth}{Formulaire d'enregistrement d'un nouvel utilisateur}{fig-superadmin-register-user}{%
L'ajout d'un utilisateur se fait à travers une fenêtre modale. Le formulaire demande les informations d'identité, l'adresse e-mail, un mot de passe temporaire ainsi que le rôle du nouvel utilisateur. Ce choix garantit une création contrôlée des accès et facilite l'affectation des responsabilités.
}

\screenpair{superadmin-new-project-manual.png}{Création manuelle d'un projet}{fig-superadmin-new-project-manual}{superadmin-new-project-ai.png}{Génération assistée d'un projet par IA}{fig-superadmin-new-project-ai}

\noindent\textbf{Explication des Figures~\ref{fig-superadmin-new-project-manual} et~\ref{fig-superadmin-new-project-ai}.} La création manuelle permet au SuperAdmin de renseigner directement le titre, la description, l'encadrant et les stagiaires lorsque le besoin est déjà formalisé. Le second parcours propose une génération assistée : le SuperAdmin décrit le besoin sous forme de prompt, puis relit, modifie si nécessaire et confirme le brouillon avant son enregistrement définitif.
\medskip

\subsection{Interfaces de l'espace Supervisor}
L'espace Supervisor est centré sur le suivi opérationnel des projets affectés. Il permet à l'encadrant de consulter les projets sous sa responsabilité, de visualiser les tâches, d'affecter le travail aux stagiaires et de suivre les livrables en attente de revue.

\screenpage{supervisor-dashboard.png}{\textwidth}{Tableau de bord de l'espace Supervisor}{fig-supervisor-dashboard}{%
Le tableau de bord Supervisor présente une vue synthétique du projet assigné, du taux d'avancement et des tâches en cours. Il met aussi en avant les tâches en attente de revue afin d'aider l'encadrant à prioriser son intervention.
}

\screenpage{supervisor-project-view.png}{\textwidth}{Liste des projets affectés au Supervisor}{fig-supervisor-project-view}{%
La vue des projets regroupe les projets auxquels le Supervisor est affecté. Chaque carte affiche l'état du projet, une description courte et l'équipe assignée. L'encadrant peut ensuite ouvrir le tableau de suivi ou ajouter une nouvelle tâche.
}

\screenpage{supervisor-project-kanban.png}{\textwidth}{Tableau Kanban de suivi des tâches}{fig-supervisor-project-kanban}{%
Le tableau Kanban structure les tâches par statut : à faire, en cours, en revue et terminées. Cette organisation donne une vision rapide de l'avancement du projet et facilite le passage d'une tâche d'un état à un autre.
}

\screenpage{supervisor-create-task.png}{\textwidth}{Création et affectation d'une tâche}{fig-supervisor-create-task}{%
La création d'une tâche permet au Supervisor de définir le titre, la description, le stagiaire assigné et l'échéance. Cette fonctionnalité transforme les objectifs du projet en actions concrètes et suivables par l'équipe.
}

\screenpage{supervisor-edit-task.png}{\textwidth}{Modification d'une tâche existante}{fig-supervisor-edit-task}{%
Le formulaire de modification offre au Supervisor la possibilité de corriger les informations d'une tâche, de changer son statut, d'ajuster la date limite ou de réaffecter le travail à un autre stagiaire.
}

\screenpage{supervisor-team-directory.png}{\textwidth}{Annuaire des stagiaires affectés au Supervisor}{fig-supervisor-team-directory}{%
L'annuaire de l'équipe donne au Supervisor une vue directe sur les stagiaires associés à ses projets. Il complète le suivi des tâches par une vision humaine de l'équipe encadrée.
}

\subsection{Interfaces de l'espace Intern}
L'espace Intern est conçu pour donner au stagiaire une vision claire de son travail quotidien. Il regroupe les tâches affectées, les échéances proches, l'état des livrables et les échanges avec l'encadrant.

\screenpage{intern-dashboard.png}{\textwidth}{Tableau de bord de l'espace Intern}{fig-intern-dashboard}{%
Le tableau de bord Intern présente les indicateurs essentiels : nombre de tâches assignées, livrables terminés et échéances à venir. La section des tâches urgentes aide le stagiaire à identifier rapidement les actions prioritaires.
}

\screenpage{intern-kanban.png}{\textwidth}{Sprint Board de l'Intern}{fig-intern-kanban}{%
Le Sprint Board organise les tâches de l'Intern selon leur statut. Le stagiaire peut suivre l'avancement de ses livrables, joindre un lien externe lorsque son travail est disponible et consulter les commentaires liés à chaque tâche.
}

\screenpage{intern-task-discussion.png}{0.78\textwidth}{Discussion autour d'une tâche}{fig-intern-task-discussion}{%
La discussion d'une tâche centralise les commentaires entre l'Intern et le Supervisor. Elle évite la dispersion des échanges et permet de conserver l'historique des remarques directement associé au livrable concerné.
}

\section*{Conclusion}
\addcontentsline{toc}{section}{Conclusion}
La réalisation a permis de mettre en place une plateforme web adaptée au suivi des projets de stage au sein de la Digital Factory de Royal Air Maroc. L'application associe une interface React par rôle, une API Spring Boot sécurisée, une base PostgreSQL et une chaîne Jenkins qui intègre les contrôles DevSecOps. Les interfaces présentées confirment la couverture des parcours essentiels du SuperAdmin, du Supervisor et de l'Intern. Le chapitre suivant présente la conclusion générale du travail et les perspectives d'évolution de la solution.
